\documentclass[10pt,letterpaper]{article}
\usepackage[T1]{fontenc}
\usepackage{amsmath,amssymb,amsthm}
\usepackage[margin=1in]{geometry}
\usepackage[most]{tcolorbox}
\usepackage{microtype}
\usepackage[hidelinks]{hyperref}
\usepackage{needspace}
\tcbset{colback=black!4,colframe=black!75,boxrule=0.7pt,arc=2pt,
 left=10pt,right=10pt,top=7pt,bottom=7pt}
\setlength{\parindent}{0pt}
\setlength{\parskip}{5pt}
\setlength{\emergencystretch}{2em}
\title{Functional Equations with Inequalities}
\author{Adhiraj Singh Anand}
\date{June 21, 2025}
\begin{document}
\maketitle
% Theory intentionally omitted. Add it here when preparing the classes.
\section{Examples}

\par\vspace{12pt}\Needspace{5\baselineskip}
\begin{tcolorbox}[breakable,title={IMO Shortlist 2023 A2},colbacktitle=black!4,coltitle=black,fonttitle=\normalfont]
\textbf{Example 1:} Let $f:\mathbb{R}\to\mathbb{R}$ be a function such that
\[f(x+y)f(x-y)\geqslant f(x)^2-f(y)^2\]
for every $x,y\in\mathbb{R}$. Assume that the inequality is strict for some $x_0,y_0\in\mathbb{R}$.

Prove that either $f(x)\geqslant0$ for every $x\in\mathbb{R}$ or $f(x)\leqslant0$ for every $x\in\mathbb{R}$.
\end{tcolorbox}
\textit{Proof:}
\[P\left(\frac{a+b}{2},\frac{a-b}{2}\right)\implies f(a)f(b)\ge f\left(\frac{a+b}{2}\right)^2-f\left(\frac{a-b}{2}\right)^2\]Let this expression be $Q(a,b)$.

\[Q(a,-b)\implies f(a)f(-b)\ge -\left(f\left(\frac{a+b}{2}\right)^2-f\left(\frac{a-b}{2}\right)^2\right)\]Adding gives $f(a)(f(b)+f(-b))\ge 0$

Now, $Q(x_0,y_0)\implies f(x_0)(f(y_0)+f(-y_0))>0$.

If $f(x_0)>0$, $f(y_0)+f(-y_0)>0$ so $Q(a,y_0)\implies f(a)\ge 0$

If $f(x_0)<0$, $f(y_0)+f(-y_0)<0$ so $Q(a,y_0)\implies f(a)\le 0$.
\par\vspace{5pt}{\footnotesize\href{https://artofproblemsolving.com/community/c6h3359723p35075360}{AoPS solution: me}}

\par\vspace{12pt}\Needspace{5\baselineskip}
\begin{tcolorbox}[breakable,title={IMO Shortlist 2023 A4},colbacktitle=black!4,coltitle=black,fonttitle=\normalfont]
\textbf{Example 2:} Determine all functions $f:\mathbb{R}_{>0}\to\mathbb{R}_{>0}$ such that
\[x\big(f(x)+f(y)\big)\geqslant\big(f(f(x))+y\big)f(y)\]
for every $x,y\in\mathbb{R}_{>0}$.
\end{tcolorbox}
\textit{Proof:}
Let $P(x,y)$ denote the assertion.

$P(x,x)\implies x\ge f(f(x))$

Assume for the sake of contradiction $f(f(x))<x$ for some $x\in \mathbb{R_+}$. Let $x-f(f(x))=z>0$.

\[P(f^n(x),f^{n-1}(x))\implies f^{n+1}(x)-f^{n-1}(x)\ge f^{n+2}(x)-f^n(x)\]So, \[x-f(f(x))\ge f(x)-f(f(f(x)))\ge \dots \ge f^{n+1}(x)-f^{n-1}(x)\ge f^{n+2}(x)-f^n(x)\]\[\implies x-f(f(x))+f(f(x))-f(f(f(f(x))))+\dots+f^{2k-2}(x)-f^{2k}(x)=x-f^{2k}(x)\ge kz\]$x>x-f^{2k}(x)\ge kz$. However, since $z>0$, this is a contradiction for large enough $k$.

So, $f(f(x))=x\implies xf(x)\ge yf(y)$

Switching $x$ and $y$ yields $xf(x)=yf(y)$, so $xf(x)=c$ and $f(x)=\frac{c}{x}$, which works $\forall c>0$.
\par\vspace{5pt}{\footnotesize\href{https://artofproblemsolving.com/community/c6h3359742p35156934}{AoPS solution: me}}
\Needspace{14\baselineskip}
\section{Practise Problems}
Here are some practise problems, in increasing order of difficulty.
\begin{enumerate}
\setlength{\itemsep}{12pt}
\Needspace{7\baselineskip}
\item \textbf{\href{https://www.imo-official.org/problems/IMO2013SL.pdf}{IMO Shortlist 2013 A3}}\par
Let $f:\mathbb{Q}_{>0}\to\mathbb{R}$ satisfy
\[f(x)f(y)\ge f(xy),\qquad f(x+y)\ge f(x)+f(y)\]
for all positive rational $x,y$. If $f(a)=a$ for some rational $a>1$, prove that $f(x)=x$ for all positive rational $x$.
\Needspace{7\baselineskip}
\item \textbf{\href{https://www.imo-official.org/problems/IMO2022SL.pdf}{IMO Shortlist 2022 A3}}\par
Find all functions $f:\mathbb{R}_{>0}\to\mathbb{R}_{>0}$ such that, for each $x>0$, there is exactly one $y>0$ satisfying
\[xf(y)+yf(x)\le2.\]
\Needspace{7\baselineskip}
\item \textbf{\href{https://www.imo-official.org/problems/IMO2011SL.pdf}{IMO Shortlist 2011 A6}}\par
Let $f:\mathbb{R}\to\mathbb{R}$ satisfy
\[f(x+y)\le yf(x)+f(f(x))\]
for all real $x,y$. Prove that $f(x)=0$ for every $x\le0$.
\Needspace{7\baselineskip}
\item \textbf{\href{https://www.imo-official.org/problems/IMO2017SL.pdf}{IMO Shortlist 2017 A8}}\par
Assume that $f:\mathbb{R}\to\mathbb{R}$ has the following property: whenever $x,y\in\mathbb{R}$ satisfy
\[(f(x)+y)(f(y)+x)>0,\]
one has $f(x)+y=f(y)+x$. Prove that
\[f(x)+y\le f(y)+x\qquad\text{whenever }x>y.\]
\end{enumerate}

\end{document}
